% ML-пайплайн обработки одной видеозаписи смены.
% Подключение в текст: % ML-пайплайн обработки одной видеозаписи смены.
% Подключение в текст: % ML-пайплайн обработки одной видеозаписи смены.
% Подключение в текст: % ML-пайплайн обработки одной видеозаписи смены.
% Подключение в текст: \input{figures/ml_pipeline}
\begin{figure}[H]
    \centering
    \begin{tikzpicture}[
        font=\small,
        node distance=18mm and 7mm,
        stage/.style={draw, rounded corners=2pt, minimum height=14mm, text width=28mm, align=center, inner sep=3pt, fill=blue!7},
        model/.style={stage, fill=orange!15},
        result/.style={stage, fill=green!10},
        config/.style={draw, dashed, rounded corners=2pt, text width=34mm, align=center, inner sep=3pt, fill=gray!8},
        flow/.style={->, >=latex, thick},
        control/.style={->, >=latex, dashed, gray!75, thick}
    ]
        \node[stage] (video) {Видеозапись смены};
        \node[stage, right=of video] (roi) {Обрезка по ROI};
        \node[model, right=of roi] (detect) {Детекция объектов};
        \node[model, right=of detect] (track) {Работа трекера};

        \node[stage, below=of track] (line) {Контроль пересечения линии};
        \node[result, left=of line] (event) {Запись о пересечении};
        \node[result, left=of event] (report) {Создание отчёта};

        \node[config, above=of detect] (project) {Конфигурация объекта};

        \draw[flow] (video) -- (roi);
        \draw[flow] (roi) -- (detect);
        \draw[flow] (detect) -- (track);
        \draw[flow] (track) -- (line);
        \draw[flow] (line) -- (event);
        \draw[flow] (event) -- (report);

        \draw[control] (project.south west) -- (roi.north);
        \draw[control] (project.south) -- (detect.north);
        \draw[control] (project.south east) -- (track.north);
        \draw[control] (project.east) -| ([xshift=8mm]line.east) -- (line.east);
    \end{tikzpicture}
    \caption{ML-пайплайн обработки видеозаписи строительного объекта.}
    \label{fig:ml-pipeline}
\end{figure}

\begin{figure}[H]
    \centering
    \begin{tikzpicture}[
        font=\small,
        node distance=18mm and 7mm,
        stage/.style={draw, rounded corners=2pt, minimum height=14mm, text width=28mm, align=center, inner sep=3pt, fill=blue!7},
        model/.style={stage, fill=orange!15},
        result/.style={stage, fill=green!10},
        config/.style={draw, dashed, rounded corners=2pt, text width=34mm, align=center, inner sep=3pt, fill=gray!8},
        flow/.style={->, >=latex, thick},
        control/.style={->, >=latex, dashed, gray!75, thick}
    ]
        \node[stage] (video) {Видеозапись смены};
        \node[stage, right=of video] (roi) {Обрезка по ROI};
        \node[model, right=of roi] (detect) {Детекция объектов};
        \node[model, right=of detect] (track) {Работа трекера};

        \node[stage, below=of track] (line) {Контроль пересечения линии};
        \node[result, left=of line] (event) {Запись о пересечении};
        \node[result, left=of event] (report) {Создание отчёта};

        \node[config, above=of detect] (project) {Конфигурация объекта};

        \draw[flow] (video) -- (roi);
        \draw[flow] (roi) -- (detect);
        \draw[flow] (detect) -- (track);
        \draw[flow] (track) -- (line);
        \draw[flow] (line) -- (event);
        \draw[flow] (event) -- (report);

        \draw[control] (project.south west) -- (roi.north);
        \draw[control] (project.south) -- (detect.north);
        \draw[control] (project.south east) -- (track.north);
        \draw[control] (project.east) -| ([xshift=8mm]line.east) -- (line.east);
    \end{tikzpicture}
    \caption{ML-пайплайн обработки видеозаписи строительного объекта.}
    \label{fig:ml-pipeline}
\end{figure}

\begin{figure}[H]
    \centering
    \begin{tikzpicture}[
        font=\small,
        node distance=18mm and 7mm,
        stage/.style={draw, rounded corners=2pt, minimum height=14mm, text width=28mm, align=center, inner sep=3pt, fill=blue!7},
        model/.style={stage, fill=orange!15},
        result/.style={stage, fill=green!10},
        config/.style={draw, dashed, rounded corners=2pt, text width=34mm, align=center, inner sep=3pt, fill=gray!8},
        flow/.style={->, >=latex, thick},
        control/.style={->, >=latex, dashed, gray!75, thick}
    ]
        \node[stage] (video) {Видеозапись смены};
        \node[stage, right=of video] (roi) {Обрезка по ROI};
        \node[model, right=of roi] (detect) {Детекция объектов};
        \node[model, right=of detect] (track) {Работа трекера};

        \node[stage, below=of track] (line) {Контроль пересечения линии};
        \node[result, left=of line] (event) {Запись о пересечении};
        \node[result, left=of event] (report) {Создание отчёта};

        \node[config, above=of detect] (project) {Конфигурация объекта};

        \draw[flow] (video) -- (roi);
        \draw[flow] (roi) -- (detect);
        \draw[flow] (detect) -- (track);
        \draw[flow] (track) -- (line);
        \draw[flow] (line) -- (event);
        \draw[flow] (event) -- (report);

        \draw[control] (project.south west) -- (roi.north);
        \draw[control] (project.south) -- (detect.north);
        \draw[control] (project.south east) -- (track.north);
        \draw[control] (project.east) -| ([xshift=8mm]line.east) -- (line.east);
    \end{tikzpicture}
    \caption{ML-пайплайн обработки видеозаписи строительного объекта.}
    \label{fig:ml-pipeline}
\end{figure}

\begin{figure}[H]
    \centering
    \begin{tikzpicture}[
        font=\small,
        node distance=18mm and 7mm,
        stage/.style={draw, rounded corners=2pt, minimum height=14mm, text width=28mm, align=center, inner sep=3pt, fill=blue!7},
        model/.style={stage, fill=orange!15},
        result/.style={stage, fill=green!10},
        config/.style={draw, dashed, rounded corners=2pt, text width=34mm, align=center, inner sep=3pt, fill=gray!8},
        flow/.style={->, >=latex, thick},
        control/.style={->, >=latex, dashed, gray!75, thick}
    ]
        \node[stage] (video) {Видеозапись смены};
        \node[stage, right=of video] (roi) {Обрезка по ROI};
        \node[model, right=of roi] (detect) {Детекция объектов};
        \node[model, right=of detect] (track) {Работа трекера};

        \node[stage, below=of track] (line) {Контроль пересечения линии};
        \node[result, left=of line] (event) {Запись о пересечении};
        \node[result, left=of event] (report) {Создание отчёта};

        \node[config, above=of detect] (project) {Конфигурация объекта};

        \draw[flow] (video) -- (roi);
        \draw[flow] (roi) -- (detect);
        \draw[flow] (detect) -- (track);
        \draw[flow] (track) -- (line);
        \draw[flow] (line) -- (event);
        \draw[flow] (event) -- (report);

        \draw[control] (project.south west) -- (roi.north);
        \draw[control] (project.south) -- (detect.north);
        \draw[control] (project.south east) -- (track.north);
        \draw[control] (project.east) -| ([xshift=8mm]line.east) -- (line.east);
    \end{tikzpicture}
    \caption{ML-пайплайн обработки видеозаписи строительного объекта.}
    \label{fig:ml-pipeline}
\end{figure}
